% Short answers to Chapter 6's exercises, lectures/L06/appendix/b_exercises.md.

\answerchapter{c6:ch}

\answer{c6:ex:1}{Sizing a counter register}
\pt{a}10 bits.
\pt{b}1023.
\pt{c}A 6-bit counter: \code{signal counter: natural range 0 to 63;}

\answer{c6:ex:2}{A counter traced by hand}
\ptnote{Trace}\begin{tabular}[t]{@{}lccccc@{}}
  Edge & 1 & 2 & 3 & 4 & 5 \\
  \midrule
  Before & 2 & 3 & 0 & 1 & 2 \\
  Captured & 3 & 0 & 1 & 2 & 3 \\
\end{tabular}
\ptnote{Why it wraps}The range needs only two bits. On the edge after \code{3} the adder produces
\code{100}; the carry has no flip-flop to go into, and the two bits that remain are \code{00}.

\answer{c6:ex:3}{Overflow in CircuitVerse}
\pt{b}One count per edge, \code{0000} up to \code{1111}, and back to \code{0000} on the sixteenth
edge. \code{carry_out} is high on one of the sixteen, while the count is \code{1111} and the
register rolls over, and nothing is connected to it. The fifth bit leaves the adder and goes
nowhere: the register has only four \code{D} inputs, so it stores the low four bits of
\code{10000}, which are \code{0000}.
\pt{c}No part of the drawing decides it: the register's four-bit width does the work, by
discarding the adder's carry-out. To count \code{0} to \code{9}, add an equality comparator that
detects \code{1001} and a clear that forces the register to \code{0000} on the next edge.
\pt{d}It wraps from \code{111} to \code{000}, after \code{7}. An N-bit counter counts \code{0} to
\code{2^N - 1} and wraps every $2^N$ edges.

\answer{c6:ex:5}{Tracing a SIPO register}
\pt{a}\begin{tabular}[t]{@{}ccc@{}}
  Edge & \code{serial_in} & \code{shift_reg} after \\
  \midrule
  1 & 1 & \code{00000001} \\
  2 & 0 & \code{00000010} \\
  3 & 1 & \code{00000101} \\
  4 & 1 & \code{00001011} \\
\end{tabular}
\pt{b}Bits 3 down to 0, holding \code{1011}. Read from bit 3 down to bit 0 they are in arrival
order: the first bit to arrive is in bit 3, the most recent in bit 0.

\answer{c6:ex:6}{Tracing a PISO register}
\pt{a}\begin{tabular}[t]{@{}lcc@{}}
  Edge & Register & \code{serial_out} \\
  \midrule
  Load & \code{1010} & 1 \\
  Shift 1 & \code{0100} & 0 \\
  Shift 2 & \code{1000} & 1 \\
  Shift 3 & \code{0000} & 0 \\
  Shift 4 & \code{0000} & 0 \\
\end{tabular}
\pt{b}Most significant bit first: bit 3, then 2, 1 and 0 of \code{parallel_in}, that is \code{1},
\code{0}, \code{1}, \code{0}.
\pt{c}\code{0000}. All four loaded bits have been shifted out of the top, and each shift filled
bit 0 with \code{serial_in = '0'}.

\answer{c6:ex:7}{SIPO and PISO in VHDL}
\pt{c}SISO: a delay line, holding a signal back a fixed number of clock cycles. SIPO: receiving a
serial bit stream, such as an incoming byte on a single wire. PISO: sending a word out over a single
wire, such as driving a chain of LEDs or an external shift-register chip. PIPO: a plain register,
holding a word.

\answer{c6:ex:8}{Accounting for the receiver}
\pt{a}\code{10110010}. The first bit to arrive, \code{1}, is in \code{data_out(7)}; the last,
\code{0}, is in \code{data_out(0)}.
\pt{b}Yes: \code{10110010}, with \code{data_ready} pulsing on the eighth bit. During the gap the
bit counter holds at \code{4}, and the register at \code{00001011}; both resume on the fifth bit. If
only the shift were gated, the counter would count the three idle edges too, so \code{data_ready}
would pulse as the fifth bit lands, with the partial byte \code{00010110} on \code{data_out}, not
pulse when the byte really completes, and stay out of step with the byte boundaries from then on.
\pt{c}\begin{itemize}
  \item Mirror-image shift: \code{data_ready} still pulses on time, but the byte-check fails,
    ``\code{serial_rx8: wrong byte on data_out for BYTE1!}'', because the byte arrives
    bit-reversed, \code{01001101}. On a real link every byte would be read reversed, with correct
    framing, so nothing but the data would look wrong.
  \item No unconditional \code{data_ready <= '0';}: \code{data_ready} rises on the first byte and
    never falls, caught on the first idle cycle by ``\code{serial_rx8: data_ready must be a
    single-cycle pulse!}''. On a real link the consumer would see a byte ready on every cycle
    after the first, and could never tell when a later byte completed.
  \item Comparison against 6: \code{data_ready} pulses on the seventh bit of the first byte,
    caught by ``\code{serial_rx8: data_ready fired before eight bits arrived!}''. On a real link the
    receiver would report 7-bit words, each missing its last bit, with the framing slipping one bit
    further from the real byte boundaries on every byte.
\end{itemize}
